% ============================================================
%  05_modele_csp.tex — Section 5 : variables, contraintes, resolution
% ============================================================
\section{Le modèle CSP : variables, contraintes, résolution}

\subsection{Variables}

Pour un benzénoïde à $h$ hexagones, on définit $h$ variables :
\[
  x_v \in \text{domaine}(v) \quad \text{pour } v = 0, 1, \ldots, h-1
\]
Le domaine est $\{6\}$ pour les hexagones gelés et $\{5, 6, 7\}$ pour les libres.

\begin{lstlisting}
x = VarArray(size=h, dom=lambda i: domains[i])
\end{lstlisting}

\subsection{Contrainte C1 — Conservation du carbone}

La formule d'Euler pour les graphes planaires impose que le nombre total de carbones soit
conservé. Cela se traduit par :
\[
  \sum_{v=0}^{h-1} x_v = 6h
\]
Autrement dit, le nombre de pentagones ($n_5$) doit être égal au nombre d'heptagones ($n_7$).

\begin{lstlisting}
satisfy(
    Sum(x) == 6 * h
)
\end{lstlisting}

\subsection{Contrainte C3 — Admissibilité du voisinage}

Pour chaque hexagone libre $v$ de degré $\geq 2$, la configuration locale
$(x_v, x_{u_1}, x_{u_2}, \ldots)$ — où $u_1, u_2, \ldots$ sont les voisins de $v$ dans le
dual — doit apparaître dans la table filtrée. C'est une \textbf{contrainte extensionnelle}.

\begin{lstlisting}
for v in free:
    if v in tables and tables[v]:
        neighbors_v = graph.neighbors(v)
        scope = [x[v]] + [x[u] for u in neighbors_v]
        satisfy(
            scope in tables[v]
        )
\end{lstlisting}

La table filtrée pour $v$ est une liste de tuples compatibles. Par exemple, pour un hexagone
libre de degré 2, les tuples sont de la forme $(x_v, x_{u_1}, x_{u_2})$.

\begin{remarque}
  La contrainte C3 peut être désactivée avec le drapeau \texttt{--no-table}. Sans elle,
  seules C1 (conservation carbone) et C4 (symétrie) restreignent les solutions, ce qui
  permet de mesurer l'impact du filtrage par la table.
  \begin{lstlisting}[style=benzai]
python main.py data/fichier.graph              # C3 active (defaut)
python main.py data/fichier.graph --no-table   # C3 desactive
  \end{lstlisting}
\end{remarque}

\subsection{Contrainte C4 — Brisure de symétrie (lex-leader)}

Deux solutions qui ne diffèrent que par un automorphisme du graphe dual décrivent la
\textbf{même molécule}. Pour éviter les doublons, on impose que chaque solution soit
lexicographiquement inférieure ou égale à toutes ses images par les automorphismes.

Pour chaque générateur $\pi$ du groupe d'automorphismes :
\[
  (x_0, x_1, \ldots, x_{h-1}) \leq_{\text{lex}} (x_{\pi(0)}, x_{\pi(1)}, \ldots, x_{\pi(h-1)})
\]

\begin{lstlisting}
for gen in generators:
    permuted = [x[gen[i]] for i in range(h)]
    satisfy(
        LexIncreasing(x, permuted)
    )
\end{lstlisting}

\begin{remarque}
  La justification théorique de cette contrainte repose sur le théorème de Whitney
  (unicité du plongement planaire pour les graphes 3-connexes) et la 2-connexité des
  benzénoïdes. Voir le document de modélisation pour la preuve complète.
\end{remarque}

\subsection{Contrainte C5 — Adjacence 5--7 (optionnelle)}

Contrainte expérimentale activée par le drapeau \texttt{--adj-57}. Elle impose que chaque
pentagone ait au moins un voisin heptagone, et inversement :
\[
  x_v = 5 \;\Longrightarrow\; \exists\, u \in \mathcal{N}(v),\; x_u = 7
  \qquad\text{et}\qquad
  x_v = 7 \;\Longrightarrow\; \exists\, u \in \mathcal{N}(v),\; x_u = 5
\]

L'implémentation utilise les constructions \texttt{If}/\texttt{Sum} de PyCSP3 :

\begin{lstlisting}
if adj_57:
    for v in range(h):
        neighbors_v = graph.neighbors(v)
        if neighbors_v:
            satisfy(
                If(x[v] == 5, Then=Sum(x[u] == 7 for u in neighbors_v) >= 1)
            )
            satisfy(
                If(x[v] == 7, Then=Sum(x[u] == 5 for u in neighbors_v) >= 1)
            )
\end{lstlisting}

\begin{remarque}
  Cette contrainte réduit l'espace de solutions aux configurations où les pentagones et
  heptagones forment des paires adjacentes (motif azulénique). Elle n'interdit pas les
  paires 5--5 ou 7--7, tant qu'il existe aussi un voisin du type complémentaire.
\end{remarque}

\subsection{Compilation et résolution}

Le modèle est compilé en XCSP3 (fichier XML), puis ACE est appelé en subprocess :

\begin{lstlisting}
# Generation du fichier XCSP3
compile(filename="model.xml")

# Appel d'ACE (solveur Java)
ace_jar = _find_ace_jar()  # dans l'installation de pycsp3
cmd = ["java", "-jar", ace_jar, "model.xml", "-s=all", "-xe"]
result = subprocess.run(cmd, capture_output=True, text=True)
\end{lstlisting}

Les drapeaux :
\begin{itemize}
  \item \texttt{-s=all} : énumérer toutes les solutions (pas seulement la première).
  \item \texttt{-xe} : exporter chaque solution au format XML dans la sortie standard.
\end{itemize}

\subsection{Parsing de la sortie ACE}

ACE écrit chaque solution sous forme de balise XML \texttt{<values>}. Le format utilise
une notation compacte : \texttt{6x3} signifie ``6 répété 3 fois''.

\begin{lstlisting}
def _parse_ace_output(output, h):
    solutions = []
    for match in re.finditer(r'<values>\s*(.+?)\s*</values>', output):
        vals_str = match.group(1).strip()
        values = []
        for token in vals_str.split():
            if "x" in token:
                val, count = token.split("x")
                values.extend([int(val)] * int(count))
            else:
                values.append(int(token))
        if len(values) == h:
            solutions.append({i: values[i] for i in range(h)})
    return solutions
\end{lstlisting}
